\section{Semantics}

\subsection{Tick Semantics (\texttt{nema.tick.v0.1})}

NEMA v0.1 defines a discrete-time, deterministic operational semantics over a fixed neural graph, identified in the spec as \texttt{nema.tick.v0.1} (spec.md:141). The tick is the single state-transition step that all backends (software reference, generated C++, and hardware flows) are required to implement consistently under the fixed-point contract (\S3.2) and determinism requirements (\S3.4).

\textbf{State.} The spec fixes the observable state used for conformance to per-node voltage and activation in signed fixed-point: voltage \texttt{V} is signed \texttt{Q8.8} stored as raw \texttt{int16}, and activation \texttt{A} is signed \texttt{Q8.8} (spec.md:143--145). The tick operates over a graph-level timestep \texttt{dt} (optional; defaults to \texttt{1.0}) and a membrane time constant \texttt{tauM} that must be positive (node-level if present, else graph-level default) (spec.md:133--134).

\textbf{Operational structure and snapshot rule.} The tick semantics is structured around a snapshot-style update discipline (as referenced by the spec's ``snapshot rule'') to ensure deterministic results independent of update ordering (spec.md:169--171). Concretely, each tick is defined as a pure function from the prior tick's state snapshot to the next tick's state: all reads of per-node state that influence the tick are conceptually taken from a single consistent snapshot, and the next-state values are written only after the tick's computations are complete.

The spec explicitly requires that simulation node iteration is index-ordered (spec.md:170). Combined with the snapshot rule (spec.md:171), this yields a deterministic evaluation for each tick even in the presence of graph structures where naive in-place updates could introduce order dependence.

\textbf{Normative reference.} The reference implementation for the tick semantics is \texttt{nema/sim.py} (spec.md:173). The spec additionally states a maintenance constraint: if the implementation and the spec disagree, both must be updated together along with regression tests in the same change (spec.md:9). This makes the semantics executable and keeps the operational definition aligned with the conformance tests that underpin the artifact.

\subsection{Fixed-Point Contract (RNE + SATURATE)}

NEMA v0.1 fixes the numeric model to a restricted, deterministic fixed-point contract. The spec states that overflow behavior is \texttt{SATURATE} only and rounding behavior is \texttt{RNE} (round-to-nearest, ties-to-even) only (spec.md:23--27). These are not configurable in v0.1 and form part of the language's semantic contract.

This explicit choice avoids floating-point execution variance that has been widely reported in both HPC and ML contexts \cite{ShanmugaveluArxiv2408.05148,AhrensTOMS2020ReproSum,BlanchardSIAM2020FABsum} and is intentionally stricter than generic IEEE-754 portability expectations \cite{IEEE754_2019}.

\textbf{Raw domains.} For a signed fixed-point value with total bitwidth \texttt{T}, the raw integer range is bounded by (spec.md:31--33):
\begin{itemize}
\item raw min: \texttt{-(2\^{}(T-1))}
\item raw max: \texttt{(2\^{}(T-1)) - 1}
\end{itemize}

This raw-domain definition is used to define saturation and to constrain behaviors at representational limits. The contract additionally constrains several primitive operations to be deterministic in edge cases, including shifts (logical right shift is defined as a bit-pattern shift), \texttt{abs} (saturates when applied to the signed minimum), and comparisons/mux/clip (deterministic with explicit ordering) (spec.md:56--58).

\textbf{Normative reference.} The spec identifies \texttt{nema/fixed.py} as the reference implementation of the fixed-point behavior (spec.md:60). In v0.1, conformance to the tick semantics is defined relative to this fixed-point contract: any backend that claims to implement \texttt{nema.tick.v0.1} must implement its arithmetic such that all observable per-tick outputs (as defined by \S3.3) are bit-identical under the stated rounding and saturation rules.

\subsection{Bit-Exactness Definition (digests)}

NEMA v0.1 defines bit-exact equivalence between executions through a per-tick digest contract. The spec states:
\begin{itemize}
\item ``Two executions are bit-exact equal when all per-tick digests match.'' (spec.md:177)
\end{itemize}

The digest is computed deterministically as follows (spec.md:179--182):
\begin{enumerate}
\item Collect \texttt{V} in node index order.
\item Pack each raw voltage as signed int16 little-endian.
\item Compute SHA-256 over the packed byte array.
\end{enumerate}

This digest definition makes the correctness criterion explicit and machine-checkable. It also makes the comparison backend-agnostic: as long as an execution can produce the same ordered sequence of raw \texttt{int16} voltages per tick, it can be validated against another execution (e.g., reference vs generated code, or software vs hardware flow) without relying on floating-point tolerances or implementation-specific traces.

Notably, the spec's digest definition is stated over \texttt{V} (spec.md:180--182). In v0.1, this makes voltage the primary observable for bit-exact equivalence, with activation \texttt{A} treated as internal to the tick computation but constrained indirectly through the same fixed-point and determinism requirements.

\subsection{Determinism Threat Model and Mitigations}

NEMA's determinism goal in v0.1 is narrow and explicit: eliminate nondeterminism in the semantics and in the conformance checks that validate backends against the reference.

\textbf{Threat model.} The primary threats addressed by the v0.1 contract are:
\begin{itemize}
\item Numeric nondeterminism from floating-point behavior differences across compilers, platforms, or optimization levels.
\item Order-dependent updates in graph computations (e.g., in-place state updates that depend on visitation order).
\item Implementation drift between the executable reference and the written specification.
\item Silent input drift from external graph/bundle references changing without detection.
\end{itemize}

\textbf{Mitigations (spec-driven).} The spec provides concrete mitigations that are required for conformance:
\begin{enumerate}
\item \textbf{Fixed-point only + fixed rounding/overflow.} By mandating \texttt{SATURATE} overflow and \texttt{RNE} rounding (spec.md:23--27), and by specifying deterministic behaviors for edge-case primitives (spec.md:56--58), NEMA removes floating-point variability from the semantic core.
\item \textbf{Canonical iteration and snapshot rule.} The semantics requires index-ordered node iteration (spec.md:170) and uses a snapshot rule to guarantee order-independent results with respect to update ordering (spec.md:171). This directly targets graph-order nondeterminism.
\item \textbf{Executable reference implementation.} The tick semantics is tied to a reference implementation (\texttt{nema/sim.py}) (spec.md:173), and the spec mandates synchronized updates to implementation + spec + regression tests when inconsistencies are discovered (spec.md:9). This is a process-level mitigation against spec/implementation divergence that would otherwise undermine reproducibility.
\item \textbf{Integrity checks for external references.} For external graph references, the spec requires that referenced paths exist and that, when a \texttt{sha256} value is not a placeholder token, it must match the file digest (spec.md:136--139). This provides a deterministic input identity and prevents ``same IR, different underlying bundle'' situations from passing silently.
\end{enumerate}

Together, these measures define a concrete determinism envelope: within the v0.1 semantics (\texttt{nema.tick.v0.1}) and fixed-point contract, executions are considered equivalent precisely when the per-tick SHA-256 digests match under the specified byte-level packing (spec.md:177--182).
